\documentclass[12pt, a4paper, twoside, english, brazil]{article}
\usepackage[top=2cm, bottom=2cm, left=2cm, right=2cm]{geometry}

\usepackage{mathtools}
\usepackage{amsfonts}
\usepackage{amssymb}
\usepackage{amsmath}
\usepackage[utf8]{inputenc}
\usepackage[portuguese]{babel}
\usepackage[T1]{fontenc}

\usepackage{enumerate}
\usepackage{multicol}
\usepackage{tikz}
\usepackage{pgfplots}

\newenvironment{solution}[1][]{\noindent\textbf{Solução#1:} }

% -----------------------------------------------------------------------------

\title{Matemática Discreta e Lógica Matemática - Lista 02}
\author{Funções de Primeiro Grau (13/04/2026)}
\date{}

\begin{document}

\maketitle

\begin{enumerate}[\bfseries 1.]
    % Questão 1
  \item Seja a função $\text{salário}(v)=1621+0.06v$, que calcula o
    salário $s$ (R\$) comissionado, em função do valor de venda $v$
    (R\$) mensal. Responda:
  \begin{enumerate}[\bfseries a)]
    \item Caso o empregado não venda no mês, qual seu salário?
    \item E se vender R\$ 150.000,00?
    \item Haverá um valor máximo de salário? Por que?
  \end{enumerate}

  \begin{solution}
  \begin{enumerate}[\bfseries a)]
    \item A venda seria $v=0$, então: $s(0)=1621+0.06\cdot0=\text R\$ 1621,00$
    \item $s(150.000)=1621+0.06\cdot 150.000=\text R\$10.621,00$
    \item Não, pois ele se comporta como uma função afim crescente,
      que aumenta infinitamente
  \end{enumerate}
\end{solution}

% Questão 2
\item Esboce o gráfico das funções $p(t)=4-4t$ e $q(t)=-8t+4$

\begin{solution}\\
  \begin{tikzpicture}
    \begin{axis}[
        axis lines = center,
        xlabel = \(x\),ylabel ={\(y\)},
        grid,
      ]

      \addplot [domain=-2:2,color=red]{-4 * x + 4};
      \addplot [domain=-2:2,color=blue]{-8 * x + 4};

      \addlegendentry{$p(x)$}
      \addlegendentry{$q(x)$}
    \end{axis}
  \end{tikzpicture}
\end{solution}

% Questão 3
\item Segundo da prefeitura de Santarém, o nível do Rio Tapajós em
2024 (24) foi de 576cm e em 2025 (25) de 710cm. Supondo que o nível
varie linearmente, encontre o modelo (função) $\text{nivel}(t)=at+b$.
Calcule o nível do rio no ano atual.

\begin{solution}
  Temos os pontos $A(24, 576)$ e $B(25, 710)$. Aplicando na equação da reta:
  \hbox

  $$
  \begin{cases}
    24a+b=576\\
    25a+b=710
  \end{cases}
  \implies
  \begin{cases}
    \text{II} -\text{I}\implies a=134\\
    \text{Aplicando $a$ em I}\implies 24\cdot 134+b=576\implies $b = -2640$
  \end{cases}
  $$

  Portanto $n(t)=134t-2640$ e $n(26)=134\cdot 26-2640=844\text{ cm}$
\end{solution}

\newpage

% Questão 4
\item Encontre a raiz da função $m(x)=-4x+8$ e desenhe o gráfico se
baseando no coeficiente linear e sua raiz.

\begin{solution}
  Para $m(x)=0$ temos $0=-4x+8\implies x=\frac{-8}{-4}=2$

  \begin{tikzpicture}
    \begin{axis}[
        axis lines = center,
        xlabel = \(x\),ylabel ={\(y\)},
        grid,
      ]

      \addplot [domain=-2:3,color=red]{-4 * x + 8};
    \end{axis}
  \end{tikzpicture}
\end{solution}

% Questão 5
\item José verificou que no 15° dia do ano o saldo em sua conta
corrente era de R\$2000,00. Devido a dívidas, no 85° seu saldo era de
$\text R\$-1500,00$. Supondo uma evolução linear de seu saldo, responda:
\begin{enumerate}[\bfseries a)]
\item Encontre o modelo $\text{saldo}(t)=at+b$
\item Quanto José tinha em conta no dia 5°
\item Em que dia o valor do saldo foi 0 (zero)?
\item Para quais intervalos de tempo $t$, o $\text{saldo}(t)>0$ e
  $\text{saldo}(t)<0.$
\end{enumerate}

\begin{solution}
\begin{enumerate}[\bfseries a)]
\item Temos os pontos $A(15, 2000)$ e $B(85, -1500)$. Portanto:\\
  \hbox

  $$
  \begin{cases}
    15a+b=2000\\
    85a+b=-1500
  \end{cases}
  \implies
  \begin{cases}
    \text{II}-\text{I}\implies 70a=-3500\implies a=-50\\
    \text{Aplicando em I}\implies 15\cdot -50+b=2000\implies b=2750
  \end{cases}
  $$

  Logo $s(t)=-50t+2750$

\item $s(5)=\text R\$2500,00$

\item $s(t)=0\implies 0=-50t+2750\implies t=55$

\item Considerando a raiz $t=55$ e sabendo que a função é decrescente:\\
  O saldo será positivo quando $t<55$\\
  O saldo será negativo quando $t>55$

\end{enumerate}

\end{solution}

\end{enumerate}
\end{document}
